\chapter{Discussion, limitations, and conclusion}\label{sec:conclusion}
A fixed local toppling rule produces broad avalanche responses as inputs encounter different configurations. The baseline height frequencies explain why the measured mean exceeds the equal-frequency estimate of $3/2$ and approaches the infinite-volume reference as the lattice grows. Mean topplings agree with the finite-grid benchmarks. Size, area and reach are strongly associated, while the source-distance comparison shows how access to a sink affects loss.

Central bias and cylindrical boundaries both increase mean topplings, through different mechanisms. Bias changes the weights assigned to input sites under an invariant recurrent measure. The independent confirmation found height differences compatible with zero. The cylinder changes the graph itself, producing higher mean heights and a response predicted by its analytic benchmark. In Eq.~\eqref{eq:benchmark}, these changes enter through $q$ and $K$, respectively. Four-grain batches instead combine successive, potentially overlapping responses. Their exact agreement with grouped sequential totals shows how event distributions can change while per-grain topplings and loss remain fixed.

The intervals reflect variation across five independent runs per primary condition, and block estimates assess temporal dependence within runs. The burn-in assessment is based on the measured observables and leaves convergence of the full configuration distribution unresolved. Fitted exponents depend on the chosen ranges, and systematic discrepancies remain even for cutoff models.

Across the comparisons, conservation fixes long-run mean loss at one per input grain. Larger lattices carry this flux through rarer events with greater loss, and batch grouping raises loss per event in proportion to its input. The combination of exact balance relations and event distributions explains how similar average transport can coexist with markedly different avalanche statistics.
